\documentclass{article}
\usepackage{graphicx} % Required for inserting images

\title{SageMath Worksheet}
\author{Rosa Orellana}
\date{August 2026}

\begin{document}

\maketitle

\begin{enumerate}
    \item For every partition $n$, $n\leq 6$, compute the change of basis matrix from the $m$-basis to the $s$-basis.  
    \begin{enumerate}
    \item What ordering is Sage using? 
    \item Is the matrix triangular? 
    \item What is the determinant? 
    \end{enumerate}
    \item Choose several partitions $\lambda$ and compute $s[\lambda].expand(4)$ for each. Make sure your examples include partitions with more and fewer than 4 parts.
    \begin{enumerate}
        \item What does this function do? 
        \item In each of your examples, use these expressions to write $s[\lambda].expand(4)$ as a sum of monomial symmetric functions. 
        \item Look for patterns in the coefficients. What do you notice? 
    \end{enumerate}
    \item For $n\leq 7$ vertices, find the chromatic symmetric function of the following class of graphs: paths, cycles, and complete graphs. 
    \begin{enumerate}
    \item For each class, expand in the monomial, elementary, power, and Schur bases.  
    \item Which basis has the simplest expression for each class of graphs? 
    \item Does the chromatic symmetric function have all positive coefficients in any of the bases? 
    \end{enumerate}

    \item Compute the chromatic symmetric function for the star graph with $n$ vertices, $St_n$ , for $1\leq n \leq 7$. Check the proposition about the expression $\mathbf{X}_{St_n}$ in the power symmetric function. 
\item Generate all trees with $n = 7$ vertices and compute their chromatic symmetric functions.  Are there two trees with the same function? Repeat with $n=8$.

\item Expand the chromatic symmetric function of each tree on $n=8$ vertices and write it in the elementary basis.  Does the coefficient of $e_{4,2,1,1}$ distinguish these trees? 

\item For every tree with 8 vertices, compute the following: 
\begin{enumerate}
\item degree sequence. 
\item diameter. 
\item number of leaves. 
\item the coefficient of $e_{n-1,1}$ and $p_{n-1,1}$.
\item Are there any correlations between the coefficients and items (a) - (c)? 
\end{enumerate}
\end{enumerate}
\centerline{\Large Mini-Projects}
\begin{enumerate}
\item[A.]
{\bf Question:} Which basis best encode trees? 

{\bf Objective:} Investigate how the chromatic symmetric function behaves in each of the bases. 
\begin{enumerate}
\item Compute the chromatic symmetric function for each tree with $n=8$ vertices. 
\item For each tree, write its chromatic symmetric function in the monomial, elementary, power, and Schur basis. 
\item Store the number of nonzero coefficients, largest coefficient, smallest coefficient. 
\item Create a table such that list every tree and number of terms in each of the basis.
\item Which basis gives the sparsest expansion? 
\item Is there a basis with all positive coefficients? 
\item Is there anything that you can tell about the tree from the coefficients? 
\end{enumerate}
\item[B.] {\bf Question: } What is the smallest and largest partitions in lexicographic order that appears with nonzero coefficient? 

{\bf Objective: } See if these partitions or their coefficients give information about the trees. 
\begin{enumerate}
\item Compute the chromatic symmetric function for each tree with $n=8$ vertices. 
\item Expand the chromatic symmetric functions in the monomial basis. 
\item For each tree write the smallest and largest partition in lexicographic largest and smallest partition. 
\item Does it encode any information about the tree? 
\item How does the coefficient change when you add a leaf to your tree? 
\item Repeat the whole exercise for the other bases of symmetric functions. 
\end{enumerate}
\end{enumerate}
\newpage
\section*{Español}
\begin{enumerate}
    \item Para cada partición de $n$, con $n\leq 6$, calcule la matriz de cambio de base de la base $m$ a la base $s$.
    \begin{enumerate}
    \item ¿Qué orden está usando Sage?
    \item La matriz que resulta, ¿es triangular?
    \item ¿Cuál es su determinante?
    \end{enumerate}
    \item Elija varias particiones $\lambda$ y calcule \texttt{s}$[\lambda]$.\texttt{expand(4)} para cada una. Asegúrese de que sus ejemplos incluyan particiones con más y menos de 4 partes.
    \begin{enumerate}
        \item ¿Qué hace esta función?
        \item En cada uno de sus ejemplos, use estas expresiones para escribir \texttt{s}$[\lambda]$.\texttt{expand(4)} como una suma de funciones simétricas monomiales.
        \item Busque patrones en los coeficientes. ¿Qué observa?
    \end{enumerate}
    \item Para $n\leq 7$ vértices, encuentre la función simétrica cromática de las siguientes clases de grafos: caminos, ciclos y grafos completos.
    \begin{enumerate}
    \item Para cada clase, expanda en las bases monomial, elemental, de potencias y de Schur.
    \item ¿Qué base tiene la expresión más sencilla para cada clase de grafos?
    \item ¿Tiene la función simétrica cromática todos sus coeficientes positivos en alguna de las bases?
    \end{enumerate}

    \item Calcule la función simétrica cromática del grafo estrella con $n$ vértices, $St_n$, para $1\leq n \leq 7$. Compruebe la proposición sobre la expresión $\mathbf{X}_{St_n}$ en la función simétrica de potencias.
    \item Genere todos los árboles con $n = 7$ vértices y calcule sus funciones simétricas cromáticas. Existen dos árboles con la misma función? Repita con $n=8$.

    \item Expanda la función simétrica cromática de cada árbol con $n=8$ vértices y escríbala en la base elemental. ¿Distingue el coeficiente de $e_{4,2,1,1}$ a estos árboles?

    \item Para cada árbol con 8 vértices, calcule lo siguiente:
    \begin{enumerate}
    \item la sucesión de grados.
    \item el diámetro.
    \item el número de hojas.
    \item el coeficiente de $e_{n-1,1}$ y $p_{n-1,1}$.
    \item ¿Hay alguna correlación entre los coeficientes y los elementos (a)--(c)?
    \end{enumerate}
\end{enumerate}
\centerline{\Large Mini-Proyectos}
\begin{enumerate}
\item[A.]
{\bf Pregunta:} ¿Qué base codifica mejor los árboles?

{\bf Objetivo:} Investigar cómo se comporta la función simétrica cromática en cada una de las bases.
\begin{enumerate}
\item Calcule la función simétrica cromática de cada árbol con $n=8$ vértices.
\item Para cada árbol, escriba su función simétrica cromática en las bases monomial, elemental, de potencias y de Schur.
\item Guarde el número de coeficientes no nulos, el coeficiente más grande y el coeficiente más pequeño.
\item Cree una tabla que enumere cada árbol y el número de términos en cada una de las bases.
\item ¿Qué base produce la expansión más esparza (\emph{sparsest})?
\item ¿Existe una base con todos sus coeficientes positivos?
\item ¿Hay algo que se pueda decir sobre el árbol a partir de los coeficientes?
\end{enumerate}
\item[B.] {\bf Pregunta:} ¿Cuáles son las particiones más pequeña y más grande en orden lexicográfico que aparecen con coeficiente no nulo?

{\bf Objetivo:} Ver si estas particiones o sus coeficientes proporcionan información sobre los árboles.
\begin{enumerate}
\item Calcule la función simétrica cromática de cada árbol con $n=8$ vértices.
\item Expanda las funciones simétricas cromáticas en la base monomial.
\item Para cada árbol, escriba la partición más pequeña y la más grande en orden lexicográfico.
\item ¿Proporciona esto alguna información sobre el árbol?
\item ¿Cómo cambia el coeficiente cuando añade una hoja a su árbol?
\item Repita todo el ejercicio para las otras bases de funciones simétricas.
\end{enumerate}
\end{enumerate}
\end{document}
